Na camada de infraestrutura, são de autoria da equipe o \textit{firmware} do nó sensor e o do \textit{gateway}, o formato do quadro, com o cabeçalho que o define para as duas pontas, o \textit{backend} com a página de monitoramento e as ferramentas que analisam o registro serial e separam reinício do nó de perda de rádio, descritos na Subseção~\ref{subsec:implementacao_infra}. São também da equipe as cinco suítes de teste executadas em placa, os testes do \textit{backend}, o \textit{firmware} de medição de energia do enlace e os programas que produziram as tabelas e figuras da Seção~\ref{sec:resultados_infra}. Não são de autoria da equipe a biblioteca \textit{RadioLib}, que controla o transceptor, o ambiente de desenvolvimento do ESP32, e o FastAPI e o SQLite, sobre os quais o \textit{backend} foi construído.

Como resultado, a infraestrutura contribui com uma bancada de enlace LoRa ponto a ponto que registra, a cada quadro, os metadados de recepção nos mesmos campos das variáveis candidatas dos modelos de detecção, e com uma coleta de referência de operação normal, de 6.007 quadros protegidos pelo Ascon-128a. Contribui também com a medição do consumo do nó na placa, que separou a parcela do processador da parcela do rádio na energia de um envio e mostrou que, no quadro de 26 bytes, a troca entre o Ascon-128a e o AES-128-GCM não altera de forma mensurável a energia de transmitir (Subseção~\ref{subsec:energia_infra}). Contribui ainda com a recuperação da coleta a partir do registro serial, que repôs sem lacuna nem duplicata os 3.575 quadros recebidos enquanto a rede esteve fora. O código, os registros das coletas e as leituras da medição de energia estão no repositório do projeto.

